\newpage
\phantomsection
\label{prf:sum-rule}

\begin{remark*}[Return]
\hyperref[thm:sum-rule]{Return to Theorem}
\end{remark*}

\begin{theorem*}[Sum Rule]
Let $A \subseteq \mathbb{R}$, let $f,g : A \to \mathbb{R}$, and let $c \in A$ be a limit point of $A$. If $f$ and $g$ are differentiable at $c$, then $f+g$ is differentiable at $c$, and
\[
(f+g)'(c)=f'(c)+g'(c).
\]
\end{theorem*}

\begin{proof}
\textbf{Professional Standard Proof.}
Let $A \subseteq \mathbb{R}$, let $f,g : A \to \mathbb{R}$, and let $c \in A$ be a limit point of $A$. Assume that $f$ and $g$ are differentiable at $c$. Then both limits
\[
\lim_{x \to c} \frac{f(x)-f(c)}{x-c}
\quad\text{and}\quad
\lim_{x \to c} \frac{g(x)-g(c)}{x-c}
\]
exist. Using algebra of limits,
\[
\lim_{x \to c} \frac{(f+g)(x)-(f+g)(c)}{x-c}
=
\lim_{x \to c} \left( \frac{f(x)-f(c)}{x-c} + \frac{g(x)-g(c)}{x-c} \right)
=
f'(c)+g'(c).
\]
Hence $f+g$ is differentiable at $c$ and $(f+g)'(c)=f'(c)+g'(c)$.
\end{proof}

\begin{proof}
\textbf{Detailed Learning Proof.}

\textbf{Step 1.} Expand the difference quotient.

\[
\frac{(f+g)(x)-(f+g)(c)}{x-c}
=
\frac{f(x)-f(c)}{x-c}
+
\frac{g(x)-g(c)}{x-c}.
\]

\textbf{Step 2.} Use differentiability.

Since $f$ and $g$ are differentiable at $c$, each term has a limit as $x \to c$:
\[
\frac{f(x)-f(c)}{x-c} \to f'(c),
\qquad
\frac{g(x)-g(c)}{x-c} \to g'(c).
\]

\textbf{Step 3.} Apply limit laws.

\[
\lim_{x \to c} \frac{(f+g)(x)-(f+g)(c)}{x-c}
=
f'(c)+g'(c).
\]

\textbf{Step 4.} Conclude.

Thus $f+g$ is differentiable at $c$ and
\[
(f+g)'(c)=f'(c)+g'(c).
\]
\end{proof}